\bookchapter{第五十七章：进阶之路——权限、通知与现代架构}

\chapterintro{能做出 App 只是开始。这一章聊聊从“会做”到“做好”的几道坎：怎么优雅地申请权限、怎么发通知不讨人厌、怎么把架构再往上抬一层，以及——在 2026 年，怎么和 AI 搭档写安卓代码。每一条都是你在真实项目里马上会遇到的。}

\booksection{一、运行时权限：先敲门，再进门}

相机、定位、通讯录、通知……这些“敏感地带”，安卓要求 App\textbf{运行时向用户申请}（Android 6 起）。规矩是：先声明（Manifest），再申请（代码），用户拒绝要优雅降级。Compose 里用官方库 Accompanist 的\texttt{rememberPermissionState}最省心：

\begin{lstlisting}[language=Kotlin]
// Manifest 里先声明：
// <uses-permission android:name="android.permission.CAMERA" />

@Composable
fun CameraButton() {
    val permission = rememberPermissionState(
        permission = android.Manifest.permission.CAMERA
    )
    Button(onClick = {
        when {
            permission.status.isGranted -> openCamera()
            permission.status.shouldShowRationale -> showExplainDialog()
            else -> permission.launchPermissionRequest()
        }
    }) { Text("拍照") }
}
\end{lstlisting}

\begin{tcbwarn}{注意：被永久拒绝怎么办？}
用户勾了“不再询问”后，\texttt{launchPermissionRequest}不会再弹窗。此时应引导用户去系统设置里手动开启（用\texttt{Intent}跳到应用设置页），并在界面里说明“这个权限用来做什么”。态度好一点，用户开权限的意愿会高很多——\textbf{申请权限的文案，值得认真写}。

\end{tcbwarn}

\booksection{二、通知：在正确的时间说有用的话}

通知是 App 触达用户的通道，也是最容易被用户“拉黑”的功能。发通知三条纪律：\textbf{只发有用的、让用户能关、别半夜骚扰}。基础用法：

\begin{lstlisting}[language=Kotlin]
// Android 13+ 必须先申请 POST_NOTIFICATIONS 权限（见上一节）

// 1. 创建通知渠道（Android 8+ 必须）
val channel = NotificationChannel(
    "reminder", "提醒", NotificationManager.IMPORTANCE_DEFAULT
)
val manager = getSystemService(NotificationManager::class.java)
manager.createNotificationChannel(channel)

// 2. 发通知
val notification = NotificationCompat.Builder(this, "reminder")
    .setSmallIcon(R.drawable.ic_notification)
    .setContentTitle("记账提醒")
    .setContentText("今天还没记账哦")
    .setPriority(NotificationCompat.PRIORITY_DEFAULT)
    .build()
if (ActivityCompat.checkSelfPermission(this,
        Manifest.permission.POST_NOTIFICATIONS) == PackageManager.PERMISSION_GRANTED
) {
    manager.notify(1, notification)
}
\end{lstlisting}

定时任务（比如“每天 20:00 提醒记账”）用\textbf{WorkManager}——它是官方推荐的“后台调度管家”，能保证任务在合适时机执行，不用自己操心电池优化。闹钟类需求用它准没错。

\booksection{三、架构进阶：依赖注入与模块化}

项目变大后，两个“甜蜜的烦恼”会找上门：

\begin{itemize}
  \item \textbf{对象创建乱}：每个页面都在自己 new 数据库、new API——用\textbf{Hilt}（依赖注入框架）统一管理，"谁要什么，开口就有人递"；
  \item \textbf{代码堆成山}：单模块项目越改越慢——拆成\textbf{多模块}（app / core / feature-xxx），各功能独立开发、独立测试。
\end{itemize}

\begin{tcbinfo}{小贴士：别过度设计}
Hilt 和模块化都是好工具，但\textbf{单人项目先别急着上}。记账本单模块 + 手动创建依赖完全够用。等代码量超过几万行、编译开始变慢、多人协作时，再按需引入。工具是服务目标的，不是用来炫耀的——这也是咱们 AI 编程书里的“简单至上”原则。

\end{tcbinfo}

\booksection{四、2026 年的新玩法：AI 辅助安卓开发}

2026 年，AI 已经深度嵌入安卓开发流程。如果你想提速，这套组合拳很成熟：

\begin{center}
\begin{tabularx}{\linewidth}{|X|X|X|}
\hline
\textbf{环节} & \textbf{AI 怎么帮} & \textbf{你的职责} \tabularnewline
\hline
写代码 & Android Studio 内置 Gemini/第三方 AI 插件补全与生成 & 描述清楚需求，审查产出 \tabularnewline
查资料 & 把官方文档和报错信息丢给 AI 解读 & 关键结论回官方文档核实 \tabularnewline
写测试 & 让 AI 按"正常/边界/异常"三类生成用例 & 确认用例符合业务语义 \tabularnewline
改 bug & 贴堆栈 + 相关代码，AI 定位修复 & 看懂改动，跑测试验证 \tabularnewline
\hline
\end{tabularx}
\end{center}

\begin{center}\small\color{gray}AI 辅助安卓开发的四个场景\end{center}

具体方法论——怎么给 AI 下指令、怎么验收、怎么避坑——完整的 11 章内容在本书第七篇《AI编程实战篇》（第37-47章）里，强烈建议配合阅读：\textbf{安卓语法交给本篇，AI 协作的心法交给第七篇}。

\booksection{五、Compose 性能优化与调试工具}

Compose 写起来爽，但“重组过多”也会卡。三个实用技巧：

\begin{itemize}
  \item \textbf{稳定参数}：给数据类加\texttt{@Immutable}或使用\texttt{derivedStateOf}，避免无关状态变化触发大面积重组；
  \item \textbf{记得用 key}：\texttt{items(list, key = \{ it.\allowbreak id \})}让列表增删时只动该动的项——不加 key，删除一项可能整表重画；
  \item \textbf{Layout Inspector 的 Compose 模式}：实时看每个 Composable 的重组次数，谁频繁重组一目了然，照着热点优化。
\end{itemize}

另外把\texttt{Compose 编译器报告}（Metrics）打开，能生成一份“哪些组件不稳定”的报告——进阶路上，数据永远比感觉可靠。

\begin{tcbpractice}{来动手：给记账本加"三件套"}
(1) 给记账本加“每天 20:00 记账提醒”（WorkManager + 通知渠道 + Android 13 通知权限）；(2) 把数据库和 API 的创建改成 Hilt 注入；(3) 试着用 AI 工具让“月度筛选”功能提速实现，然后按三层验收法检查产出。三件套做完，你的项目就有“正规军”的样子了。

\end{tcbpractice}

\begin{tcbsummary}{本章小结}
\begin{itemize}
  \item 敏感权限：Manifest 声明 + 运行时申请 + 被拒后优雅引导
  \item 通知三纪律：只发有用的、能关、不骚扰；渠道 + 权限缺一不可
  \item 定时任务用 WorkManager，别自己造轮子
  \item Hilt 管对象、模块化管代码——单人小项目先别急着上
  \item AI 辅助四场景：写码、查资料、写测试、改 bug；方法论见第七篇（第37-47章）
\end{itemize}

\end{tcbsummary}
